\SourceLeaf{208}\begin{center}
[\textsc{Bessel} \emph{an} \textsc{Gauss}. \emph{Königsberg i. Pr., 10. Februar 1829.}]
\end{center}

\{.\ .\ .\ .\ .\ .\ Ich würde sehr bedauern, wenn Sie sich »durch das Geschrei der
Böotier« davon abhalten ließen, Ihre geometrischen Ansichten darzulegen.
Durch das, was \textsc{Lambert} gesagt hat und was \textsc{Schweikart} mündlich äußerte,
ist mir klar geworden, dass unsere Geometrie unvollständig ist und eine Korrektur erhalten sollte, die hypothetisch ist und verschwindet, wenn die Summe der
Winkel des ebenen Dreiecks $=180^{0}$ ist. Das wäre die \textls{wahre}
Geometrie, die euklidische die \textls{praktische}, wenigstens für Figuren auf der
Erde. .\ .\ .\ .\ .\ .\}

\rule{4cm}{0.5pt}

\Needspace{16\baselineskip}
\begin{center}
\textsc{Gauss} an \textsc{Bessel}. Göttingen, 9. April 1830.
\end{center}

.\ .\ .\ .\ .\ .\ .\ Die Leichtigkeit, mit der Sie auf meine Ansichten über die Geometrie
eingegangen sind, hat mir wahre Freude gemacht, zumal so wenige
dafür offen sind. Nach meiner tiefsten Überzeugung nimmt die Raumlehre in unserem Wissen a priori eine ganz andere Stellung ein als die reine
Größenlehre. Unserer Kenntnis von jener fehlt durchaus die vollständige Überzeugung von ihrer Notwendigkeit (also auch von ihrer absoluten
Wahrheit), die der letzteren eigen ist. Wir müssen in Demut zugeben,
dass, wenn die Zahl \textls{bloß} ein Produkt unseres Geistes ist, der Raum auch außerhalb
unseres Geistes eine Realität hat, der wir a priori ihre Gesetze nicht vollständig vorschreiben können. .\ .\ .\ .\ .\ .\ .\ .

\rule{4cm}{0.5pt}

